\section{Local arithmetic of the minors}\label{sec:local}
Assume simultaneous rationality and fix the positive common denominator
$Q$ in~\eqref{eq:arithmetic-form}. We normalize $v_p(p)=1$ and allow
$v_p(0)=+\infty$ in lower-bound statements. Arithmetic improvements for
hypergeometric linear forms exploit prime divisibility beyond a naive common
denominator; see \cite{Zudilin04,Viola03}. Here the relevant gain is at the
level of square minors of local moment matrices.

\subsection{A bound valid at every prime}
Put $L_p=\lfloor\log_p(800n)\rfloor$.
\begin{lemma}\label{lem:small-primes}
Every entry in the matrix defining~\eqref{eq:minor} satisfies
\begin{equation}\label{eq:small-entry}
 v_p(F_4(E_iE_j))\ge-40L_p-v_p(Q).
\end{equation}
Consequently $v_p(D_n)\ge-40KL_p-Kv_p(Q)$.
\end{lemma}
\begin{proof}
We keep the Laurent order in the estimate rather than bounding all
coefficients separately. At each pole, substitute $x=a+p^{L_p}u$.
All nonzero root differences under consideration have absolute value at
most $800n$. For each prime power $p^e\le800n$, the centered root counts
and the factorial floors in the normalization give a lower bound $-30$
for the weight contribution. Indeed a centered count of radius $b$ differs
from $2b/p^e$ by at most one, and the linear terms cancel against the
normalization floors; the ten normalization floor differences lose at most
ten and the five numerator and fifteen denominator counts lose at most
twenty. The explicit factor $x$ has nonnegative contribution. Summing over
$e$ gives the scaled Laurent bound $-30L_p$ for $w_n$.

For $i=j+1\ge1$ the exact factorization
\[
 E_i(x^2)=\frac{2xD_j(x)}{(2j+2)!}
\]
shows that each basis factor loses at most $3L_p$. At each prime power,
the centered count for $D_j$ minus the factorial floor is at least $-3$;
the factor $x$ again contributes nonnegatively. The constant basis vector
requires no loss. It follows that the coefficient of $(x-a)^{-k}$ in
$w_n(x)E_i(x^2)E_j(x^2)$ has valuation at least $(-36+k)L_p$.
The harmonic factor of order $k+4$ has valuation at least
$-(k+4)L_p$, since its summation indices are at most $800n$.
The fixed rational zeta coefficients cost at most $v_p(Q)$.
Thus each local summand has valuation at least $-40L_p-v_p(Q)$.
Summation over poles preserves the lower bound. Expanding a $K$-square
determinant gives the last assertion.
\end{proof}
This bound, including higher prime powers, will be used on the whole
omitted interval $p\le n/20$. It replaces the earlier, weaker small-prime
estimate in the initial manuscript.

\subsection{Nonzero residue classes}
Let $p>\sqrt{800n}$.  Eventually such primes do not divide $Q$.
All factorial arguments, root differences, harmonic indices, and
basis factorial arguments are then below $p^2$.  Put
\begin{align}
 \nu_p&=\sum_{t\in T}\floor{2tn/p}-10\floor{33n/p},\label{eq:nu}\\
 c_a(r)&=\floor{(a+r)/p}+\floor{(a-r)/p}+1,\notag\\
 \lambda_r&=\nu_p+5c_b(r)-\sum_{i=1}^{15}c_{a_i}(r),\label{eq:lambda}\\
 c_r&=\sum_{t\in T}c_{tn}(r),\qquad 1\le r<p/2.\label{eq:capacity}
\end{align}
The number $c_a(r)$ counts the integers in $[-a,a]$ congruent to $r$
modulo $p$.  The two $x$-classes $r,-r$ form one class for $Y=x^2$.
Group the partial-fraction expression~\eqref{eq:arithmetic-form}
accordingly, obtaining a local functional and its bilinear moment
matrix $M_r$.

\begin{lemma}\label{lem:local-entry}
The rank of $M_r$ is at most $c_r$.  In the translated monomial basis,
\begin{equation}\label{eq:entry-bad}
 v_p\bigl(M_r((Y-r^2)^i,(Y-r^2)^j)\bigr)
       \ge\lambda_r+i+j-4.
\end{equation}
Let $j_{\max}$ be the largest absolute actual pole in this class.
If $p>a_0+j_{\max}$, the right side improves to
$\lambda_r+i+j+1$.  If in addition there is only one positive
$Y$-node, every entry in an integral polynomial basis has valuation
at least $\nu_p$.
\end{lemma}
\begin{proof}
At each positive absolute pole $j$, the local expression is a linear
combination of derivatives of $P(Y)$ at $Y=j^2$, of orders below $o_j$.
The paired negative pole gives the same $Y$-node.  The sum of these
jet lengths in the class is precisely $c_r$: counting the $r$ class
in the symmetric interval already accounts for both possible signs
of the positive representative.  This proves the rank bound, without
an extra factor of two.

For the entry estimate set $x=r+pu$, or $x=-r+pu$ at the opposite
class.  Factoring out the powers of $p$ from the numerator and
denominator of $w_n$ gives $p^{\lambda_r}$ times a rational function
whose scaled Laurent coefficients at the relevant poles are integral.
Factors not in this residue class are units; distinct scaled poles
have unit differences because the unscaled differences have absolute
value less than $p^2$.  The explicit factor $x$ is a unit as well.
Moreover $x^2-r^2$ has a factor $p$ in these coordinates.
Thus for a polynomial factor $(Y-r^2)^{i+j}$, an order-$k$ Laurent
coefficient in the original $x$ coordinate has valuation at least
$\lambda_r+i+j+k$.  The bracket in~\eqref{eq:arithmetic-form}
loses at most $k+4$ powers of $p$, which proves~\eqref{eq:entry-bad}.

In a good class, namely $p>a_0+j_{\max}$, all harmonic indices are
less than $p$; the brackets are integral.  Every Laurent order is at
least one, giving the improved entry estimate.  Finally, in a good
single-node class, remove the pole factor at that node from the
already cancelled denominator.  All its remaining factors are
$p$-units at the node.  The numerator is an integral polynomial times
$\mathcal N_n$, so all its Laurent coefficients have valuation at
least $v_p(\mathcal N_n)=\nu_p$.  Integral polynomial multiplication
preserves this bound, and so do the good harmonic weights.
\end{proof}

\begin{lemma}\label{lem:local-minor}
Every $m$-minor of $M_r$, in any polynomial family with integral
monomial coefficients, has valuation at least $f_r(m)$, where
\begin{equation}\label{eq:costs}
 f_r(m)=
 \begin{cases}
 m^2+(\lambda_r-5)m,&\text{class not good},\\
 m^2+\lambda_rm,&\text{good, more than one node},\\
 \max\{m^2+\lambda_rm,\nu_pm\},&\text{good, one node}.
 \end{cases}
\end{equation}
The domain is $0\le m\le c_r$; larger minors vanish.
Each displayed integer cost function is convex.
\end{lemma}
\begin{proof}
For an $m$-minor in the translated monomial basis, selected exponents
on either side are distinct nonnegative integers, with sum at least
$m(m-1)/2$.  Summing~\eqref{eq:entry-bad} along a permutation gives
\[
 m(\lambda_r-4)+m(m-1)=m^2+(\lambda_r-5)m.
\]
The good estimate adds $5m$.  The single-node entry bound gives the
additional lower bound $m\nu_p$.  Translation is integral and
unitriangular; Cauchy--Binet transfers each lower bound to every
integral polynomial family.  The rank statement follows from
Lemma~\ref{lem:local-entry}.  Quadratic costs have increasing marginal
increments, and the maximum of the displayed quadratic and affine
costs is convex, including on their integer domains.
\end{proof}

\subsection{A finite dual bound}
The determinant of a sum of local matrices is expanded by assigning its
columns to residue classes and then applying Laplace expansion to the
corresponding groups of rows. Each term is a product of local square minors
of sizes $m_r$, with $\sum_r m_r=K$. For any integer $t$ define
\[
 H_r(t)=\max_{0\le m\le c_r}\{tm-f_r(m)\}.
\]
Then $f_r(m)\ge tm-H_r(t)$ for every admissible $m$, and hence
\begin{equation}\label{eq:dual}
 v_p(D_n)\ge tK-\sum_r H_r(t)-F_p,\qquad
 F_p=\sum_{i\in I}v_p((2i)!)+\sum_{j\in J}v_p((2j)!).
\end{equation}
The factorial cost is applied once after combining the classes. For odd
$p$ the factor two in~\eqref{eq:E} is a unit. This argument needs only a
feasible dual bound, not an optimizer or a proof of optimality.
The class $r=0$ and classes on partition boundaries are included below by
an explicit coarse penalty. Thus no residue class is silently discarded.

The local estimates are initially made for integral numerator polynomials.
The change from translated monomials to those polynomials is integral;
Cauchy--Binet preserves the lower bound. The subsequent division by the
selected row and column factorials gives exactly $F_p$ in~\eqref{eq:dual}.
All of this takes place in the physical variable $Y=x^2$, the same variable
as in the analytic determinant.

\section{Finite arithmetic cells and the prime sum}\label{sec:prime}
\subsection{Uniform rational boxes}
Use the scale
\[
 u=p/n,\qquad v=r/n,\qquad \widehat C=361/2,
\]
so the certified prime range is $1/20<u\le361$. Let
$\mathcal A=(100,67,63,60,58,57,54,52,49,47,44,42)$.
For a rational cell $(l,h]$, partition $[0,h/2]$ into finitely many
rational boxes $[a_B,b_B]$. The actual residue range $[0,u/2]$ is contained
in this interval. For each endpoint $A\in\mathcal A$, the box stores integer
bounds $L_A^\pm,U_A^\pm$ satisfying, at all four corners,
\begin{equation}\label{eq:corner}
 L_A^\pm u\le A\pm v\le(U_A^\pm+1)u,
 \qquad (u,v)\in\{l,h\}\times\{a_B,b_B\}.
\end{equation}
Each difference in these inequalities is affine in $(u,v)$, so the bounds
hold throughout the rectangle. For $a_B<v<b_B$, the nonzero coefficient
of $v$ makes the upper inequality strict. Thus
\[
 L_A^\pm\le\left\lfloor\frac{A\pm v}{u}\right\rfloor\le U_A^\pm
 \quad(l\le u\le h,\ a_B<v<b_B).
\]
This is the corner-to-interval step; a midpoint evaluation would not suffice.
Write $c_A^-=L_A^++L_A^-+1$ and $c_A^+=U_A^++U_A^-+1$.
Normalization floors give an integer $\nu^-\le\nu_p$ valid on the cell.
The box stores integers $c_B,\lambda_B,\nu_B$ such that
\begin{align*}
 c_B&\ge\sum_{A\in T}c_A^+,&
 \lambda_B&\le\nu^-+5c_{100}^- -5c_{67}^+-\sum_{A\in T}c_A^+,
 &\nu_B&\le\nu^-.
\end{align*}
A good-class flag is permitted only when the corresponding pole-location
inequalities hold at every corner. They imply that all harmonic indices
are less than $p$ for an interior residue. A single-node flag additionally
requires $U_{63}^++U_{63}^-\le0$. For a nonempty class this leaves only one
node. These sufficient tests may sacrifice sharpness but never presume
that a classification at one sample point persists across the cell.

Let $g_B(m)$ be the cost~\eqref{eq:costs} formed with the box's lower
parameters and its certified flags. Choose a cell price $t\in\Z$ with
$|t|\le20000$ and nonnegative integers $H_B$ satisfying
\begin{equation}\label{eq:box-dual}
 tm-H_B\le g_B(m)\qquad(0\le m\le c_B,\ m\in\Z).
\end{equation}
A quadratic cost $m^2+am$ reduces this check to the endpoints and integers
nearest $(t-a)/2$. For a cost $\max(m^2+am,\nu m)$, split at
$m=\nu-a$ and check the kink as well. On each resulting interval the
function $tm-g_B(m)$ is either affine or a concave quadratic. This proves
that the finite candidate checks used by the verifier control every
integer $m$ in the stated domain.

\subsection{Boundary residues and the factorial staircase}
There are at most one more boundary points than boxes. For $u\ge1/20$,
every centered count needed here is at most $4002$, every residue capacity
is at most $40020$, and the local weight is at least $-100000$.
Consequently the cost is bounded below by $m^2-100005m$, and for
$|t|\le20000$, $0\le m\le40020$,
\[
 tm-10^{12}\le m^2-100005m.
\]
The fixed penalty $10^{12}$ therefore covers each boundary class, including
$r=0$. This deliberately generous constant affects only an error of lower
asymptotic order.

The number of residues in an interior box is at most
$n(b_B-a_B)+1$. Since $H_B\ge0$, summing the box penalties gives
\[
 \sum_r H_r(t)\le n\sum_B(b_B-a_B)H_B+R_c,\qquad
 R_c=10^{12}(J_c+1)+\sum_BH_B,
\]
where $J_c$ is the number of boxes in the cell. Put $k=\lfloor361/l\rfloor$.
For $u>l$, and $800n<p^2$, the full factorial cost satisfies
\begin{equation}\label{eq:factorial-cost}
 F_p\le n\left(361k-\frac{lk(k+1)}2\right)+2k.
\end{equation}
To see this, enlarge the selected index sets to all $0\le i<N$ and
write $\lfloor2i/p\rfloor$ as the number of positive integers $j$ with
$jp\le2i$. There are at most $k$ levels; the number of indices at level
$j$ is at most $(361n-jp)/2+1$. Sum these counts on both sides of the
minor and use $p>ln$.

If the cell lower value $L_c\in\Q$ satisfies
\begin{equation}\label{eq:cell-lower}
 L_c\le t\widehat C-\sum_B(b_B-a_B)H_B
                   -361k+\frac{lk(k+1)}2,
\end{equation}
then $|K-\widehat Cn|\le4$, together with~\eqref{eq:dual}, gives
\begin{equation}\label{eq:cell-valuation}
 v_p(D_n)\ge nL_c-E_c,\qquad E_c=R_c+4|t|+2k.
\end{equation}
This holds whenever $l<p/n\le h$, $800n<p^2$, $p\nmid Q$, and the
minor is nonzero. In particular it holds throughout the certified prime
range for $n\ge320000$ and $n\ge20Q$.

\subsection{The exact arithmetic certificate}
The data file \path{numerical/global28/arithmetic_cells.txt} gives
$171{,}504$ adjacent rational cells covering $(1/20,361]$, together with
their integer prices and claimed rational lower values. The verifier
reconstructs $3{,}810{,}331$ boxes, including zero-capacity boxes, and
checks~\eqref{eq:corner}--\eqref{eq:cell-lower} in exact rational arithmetic.
Each cell-integral contribution is rounded downward at scale $10^{-8}$.
The verified sum is bounded below by
\begin{equation}\label{eq:padic-cert}
 \sum_c(h_c-l_c)L_c\ge
 -\frac{1434459186663}{12500000}
 =-114756.73493304>-114800.
\end{equation}
The Lean finite predicate requires only the weaker last lower bound
$-114800$. It also requires the exact cell count, ordered endpoints,
adjacency, finite box inequalities, and price bounds.
Its validity is the single externally checked arithmetic-grid admission.
The transfer from that finite predicate to~\eqref{eq:cell-valuation} is
proved in Lean.

\subsection{Chebyshev functions and omitted primes}
Define
\[
 \theta(x)=\sum_{p\le x}\log p,\qquad
 \psi(x)=\sum_{p^a\le x}\log p.
\]
The prime number theorem gives $\theta(x)\sim x$ and $\psi(x)\sim x$;
the formal proof uses the pinned development \cite{PNT}. For each fixed
cell,
\[
 \frac{\theta(h_cn)-\theta(l_cn)}n\longrightarrow h_c-l_c.
\]
Only finitely many such limits are needed. This avoids a limiting
optimization problem or an exchange of a prime sum with an unknown
continuous minimizer.

For $T\le X$, expanding $\psi$ by prime powers gives
\begin{equation}\label{eq:omitted-primes}
 \sum_{p\le T}\lfloor\log_pX\rfloor\log p
 \le\psi(X)-\theta(X)+\theta(T).
\end{equation}
Indeed the first power contributes $\theta(T)$, and the higher powers
form a sub-sum of $\psi(X)-\theta(X)$. Lemma~\ref{lem:small-primes}
therefore bounds the omitted-prime loss by
\[
 40K\{\psi(800n)-\theta(800n)+\theta(n/20)\}.
\]
After division by $(100n)^2$, this tends to
\begin{equation}\label{eq:small-allowance}
 \frac{40(361/2)}{10000\cdot20}=\frac{361}{10000}=0.0361.
\end{equation}
The higher powers contribute zero to this normalized limit. Above $361n$
all relevant factorial arguments, root differences and harmonic indices
are less than $p$; hence the determinant is $p$-integral apart from the
fixed denominator $Q$. Eventually these primes do not divide $Q$ either.
Over all primes the denominator cost is at most $K\log Q=O(n)$.

\subsection{The eventual global lower bound}
Let $E_G=1+\sum_c|E_c|$, a fixed constant. Combining the preceding
bounds with the rational product formula yields
\begin{align}
 \log|D_n|\ge{}&n\sum_c L_c\{\theta(h_cn)-\theta(l_cn)\}
                   -E_G\theta(361n)\notag\\
 &-40K\{\psi(800n)-\theta(800n)+\theta(n/20)\}
                   -K\log Q.\label{eq:global-log-bound}
\end{align}
The cell-boundary convention is $(l_c,h_c]$, so each prime is counted
exactly once. The fixed error and fixed-denominator terms are $o(n^2)$.
Thus the finite predicate gives
\begin{equation}\label{eq:prime-liminf}
 \liminf_{n\to\infty}\frac{\log|D_n|}{(100n)^2}
 \ge\frac{-114800-361}{10000}=-11.5161.
\end{equation}
The actual checked sum in~\eqref{eq:padic-cert} gives the stronger numerical
value $-1438971686663/125000000000=-11.511773493304$ after the same allowance.
For the final theorem it suffices to use the strict slack in
$-11.5161>-11.52$. We obtain, for every sufficiently large $n$,
\begin{equation}\label{eq:P-certified}
 \log|D_n|\ge P_*(100n)^2,\qquad P_*=-\frac{288}{25}=-11.52.
\end{equation}
The cutoff may depend on $Q$ and the fixed certificate constants; no
numerical value for it is asserted. This is the arithmetic rate used in
the completed proof. The stronger midpoint-allocation rate printed in
the 12 October manuscript is superseded and is not used here.
